\documentclass[11pt]{article}
\usepackage[margin=1in]{geometry}
\usepackage{amsmath,amssymb,amsthm,booktabs,microtype}
\usepackage[T1]{fontenc}
\usepackage{lmodern}
\setlength{\emergencystretch}{1em}
\usepackage[colorlinks=true,linkcolor=blue,citecolor=blue,urlcolor=blue]{hyperref}
\newtheorem{theorem}{Theorem}
\newtheorem{lemma}[theorem]{Lemma}
\newtheorem{proposition}[theorem]{Proposition}
\newcommand{\F}{\mathbb F}
\newcommand{\Tr}{\operatorname{Tr}}
\newcommand{\im}{\operatorname{im}}
\newcommand{\Span}{\operatorname{span}}
\newcommand{\End}{\operatorname{End}}
\title{Certified recovery of trace-determinant structure\\under low-dimensional corruption}
\author{Daniel Eric Fredriksen\\Quantyra Inc}
\date{28 October 2026 --- working manuscript, not a released version}
\begin{document}
\maketitle
\begin{abstract}
Let an explicitly represented alternating trilinear tensor on $\F_2^{3s}$ be
the sum of a nonzero trace-determinant tensor over an unknown field
$\F_{2^s}$ and an arbitrary alternating perturbation supported on at most $R$
binary covectors. We give a bounded randomized reconstruction using four
contractions, scalar-field recognition, and a two-translate defect operator.
When $s\geq2R+5$, at most $8\cdot2^{4R}h$ trials return a verified compatible
representation with probability at least $1-2^{-h}$. Every accepted result is
checked exactly, including outside the sufficient promise. Conservative bit
bounds are $\widetilde O(2^{4R}h n^{12})$ time and
$\widetilde O(n^8+\log(h+1))$ working space, where $n=3s$.
The result concerns supported algebraic corruption, not arbitrary noise or
unique identification of a planted field. The companion contains executable
finite checks. Complete Lean certification is in progress; this draft does
not claim that the main theorem or the Python implementation is formally verified.
\end{abstract}

\section{Problem and contribution}
Put $V=\F_2^n$, $n=3s$. An alternating trilinear tensor is given by the
$d=\binom n3$ coefficients $C(e_i,e_j,e_k)$ for $i<j<k$, in lexicographic
order. Trilinearity and alternation determine all other coefficients. Over
$\F_2$, alternation implies symmetry, and coefficients with a repeated index
are zero. This is a dense exact input model, of size $\Theta(n^3)$ bits;
the input is not an oracle or a succinct circuit.

For a binary subspace $U\subseteq V^*$, write $U^\perp$ for its annihilator.
We say $E$ is supported on $U$ when it factors through $V/U^\perp$ in every
argument. Thus $E(x,y,z)=e(u(x),u(y),u(z))$ for a surjective binary linear map
$u:V\to\F_2^r$ whose coordinate covectors form a basis of $U$, and an
alternating trilinear form $e$. No distribution is assumed for $U$ or $e$.
This condition is different from a bound on tensor rank, a sparse set of
corrupted entries, or independent statistical errors.

The proposed contribution is the specific worst-case supported-corruption
guarantee below, with a complete certificate and charged search. The
trace-determinant construction, adjoint equations, and cancellation of internal
perturbations have substantial precedents; Section~\ref{sec:prior} describes
the comparison. Neither priority nor a new general denoising principle is claimed.

\begin{theorem}[Bounded compatible recovery]\label{thm:main}
Let $0\leq R\leq n$, $h\geq1$ be integers. Suppose there exist a field
$K\cong\F_{2^s}$, a binary linear isomorphism $f:V\to K^3$, $c\in K^\times$,
and an alternating tensor $E$ supported on $U_*$ with $\dim_{\F_2}U_*\leq R$
such that
\begin{equation}\label{eq:model}
 C(x,y,z)=\Tr_{K/\F_2}\bigl(c\det(fx,fy,fz)\bigr)+E(x,y,z).
\end{equation}
If $s\geq2R+5$, the algorithm in Section~\ref{sec:algorithm}, given only
$(n,C,R,h)$, succeeds with probability at least $1-2^{-h}$ using independent
fair bits. It returns some verified representation of the form
\eqref{eq:model}, with nonzero coefficient and support dimension at most $R$.
It uses at most $N=8\cdot2^{4R}h$ trials and $(4n+s)N$ fair bits.
On every tape it terminates with a verified representation or explicit failure.
Every accepted representation is sound even if the sufficient promise fails.
The conservative bit-time and working-space bounds are
\begin{equation}\label{eq:cost}
 \widetilde O(2^{4R}h n^{12}),\qquad
 \widetilde O(n^8+\log(h+1)),
\end{equation}
where the soft notation hides logarithmic factors in $n$ and $h+1$.
\end{theorem}

The theorem does not assert uniqueness of $K$, $f$, $c$, or $U_*$. In
particular a valid alternative representation is a successful output.
Malformed input is rejected before any random draw. Failure on a well-formed
input is not a certificate that no compatible representation exists.

\section{Four contractions expose the scalar field}
For this section identify $V$ with $K^3$ through the unknown $f$ and write
$C_0(x,y,z)=\Tr(c\det(x,y,z))$. A useful elementary fact is that
$(x,y)\mapsto\Tr(\lambda xy)$ is a nondegenerate binary pairing for every
$\lambda\ne0$. Indeed the trace polynomial
$X+X^2+\cdots+X^{2^{s-1}}$ is nonzero of degree less than $2^s$, so the
trace functional is nonzero; multiplication by any nonzero element is bijective.

\begin{lemma}[Four-contraction identification]\label{lem:four}
Let $v_1,v_2,v_3$ be a $K$-basis and let
$v_4=a v_1+b v_2+d v_3$. If at least one of $a,b,d$ generates $K$ as an
$\F_2$-algebra, then the binary vector space
\[
 \mathcal B=\{T\in\End_{\F_2}(V):T^{\mathsf t}B_j=B_jT,\ j=1,2,3,4\},
 \qquad B_j(x,y)=C_0(x,y,v_j),
\]
is exactly the $s$-dimensional space of scalar multiplications by $K$.
\end{lemma}
\begin{proof}
The radical of $B_j$ for $j\leq3$ is $K v_j$: if $x$ is outside that line,
the map $y\mapsto\det(x,y,v_j)$ is a nonzero $K$-linear functional onto $K$,
and the trace pairing is nondegenerate. A balanced endomorphism preserves
each radical, since $B_j(Tx,y)=B_j(x,Ty)=0$ for $x$ in the radical. It is
therefore diagonal in this three-block decomposition, with binary maps
$A_1,A_2,A_3:K\to K$.

Use the pairing $\langle x,y\rangle=\Tr(\lambda xy)$, where
$\lambda=c\det(v_1,v_2,v_3)\ne0$. Balancing the three forms on their
off-diagonal blocks gives $A_2^*=A_3$, $A_1^*=A_3$, and $A_1^*=A_2$.
Consequently $A_1=A_2=A_3=A=A^*$. In the fourth form, the three
off-diagonal pairings are weighted by $a,b,d$. Their balance equations say
that $A$ commutes with multiplication by each of these coefficients. A map
commuting with a field generator commutes with all its polynomial expressions,
so is $K$-linear. On the one-dimensional $K$-space $K$, it is multiplication
by a scalar. Conversely, scalar multiplication balances all four forms by
$K$-multilinearity of the determinant. This proves equality.
\end{proof}

Three contractions alone impose self-adjointness of $A$ but not $K$-linearity.
The space of self-adjoint maps has binary dimension $s(s+1)/2$, because
composition with a fixed nondegenerate pairing identifies it with symmetric
binary matrices. Thus the fourth equation carries necessary additional information
for this argument.

\section{Recovering the defect with two translates}
For a tensor $E$ let its third mode image be
\[
 W(E)=\Span_{\F_2}\{E(x,y,\cdot):x,y\in V\}\subseteq V^*.
\]
For alternating $E$, symmetry makes this the mode image in every argument.
It is the smallest support space: if $z\in W(E)^\perp$, every $E(x,y,z)$
vanishes, and symmetry gives the corresponding statements in the other
arguments. The tensor therefore descends to $V/W(E)^\perp$ in all arguments.
Conversely every support space contains all contracted covectors.

For a candidate scalar action $T$, define
\begin{equation}\label{eq:defect}
 D_T C(x,y,z)=C(Tx,y,z)+C(x,Ty,z).
\end{equation}
If $T$ is multiplication by $t\in K$, then $D_T C_0=0$.

\begin{lemma}[Two-translate image recovery]\label{lem:two}
If $E$ is supported on $U$ and $U\cap T^*U=0$, with $T$ invertible on
$U$ in the sense that $\dim T^*U=\dim U$, then
$W(D_T E)=W(E)$. Equivalently it suffices that a basis for $U$, followed by
its $T^*$-translate, is linearly independent. The empty support case also holds.
\end{lemma}
\begin{proof}
Write $E(x,y,z)=e(u(x),u(y),u(z))$, where $u:V\to\F_2^r$ records a
basis of $U$. The joint map $x\mapsto(u(x),u(Tx))$ is onto, because its
$2r$ coordinate covectors are independent. Every contracted covector of
$D_T E$ is a sum of two contracted covectors of $E$, proving one inclusion.
For any $x,y$, choose $a$ with $u(a)=0$ and $u(Ta)=u(x)$. Then
\[
 D_T E(a,y,z)=e(u(x),u(y),u(z))+e(0,u(Ty),u(z))=E(x,y,z)
\]
for all $z$. This proves the reverse inclusion. When $r=0$ both tensors vanish.
\end{proof}

The algorithm needs only $W(D_T C)$, obtained by binary elimination on
explicit coefficient slices. It does not test cleanness against a supplied
$U_*$. Under the good event this recovered space is $W(E)\subseteq U_*$.
An alternating tensor supported on an $r$-dimensional space is a linear
combination of its $\binom r3$ exterior basis forms. This follows by choosing
dual coordinates on the quotient: alternation sets repeated indices to zero,
and the strictly increasing triples specify all remaining coefficients.

\section{Conditional sampling and amplification}
Put $q=2^s$, $r=\dim U_*$, and $A=U_*^\perp$. Four independent uniform
vectors in $V$ all lie in $A$ with probability $2^{-4r}\geq2^{-4R}$.
Conditional on this event they are independent uniform elements of $A$;
their field coordinates need not be uniform in $K$.

The first three fail to form a $K$-basis with probability at most
\[
 \frac{1+q+q^2}{|A|}=2^r\frac{1+q+q^2}{q^3}.
\]
At the successive stages the forbidden sets are the origin, a $K$-line,
and a $K$-plane. Intersecting with $A$ only decreases their sizes. Conditional
on any selected basis triple, the first coordinate of $v_4$ has fibers of
size at most $q^2$ within $A$.

Define
\[
 \tau(s)=\sum_{d\mid s,\ d<s}2^d.
\]
A non-generator has minimal polynomial degree a proper divisor of $s$ and
belongs to the corresponding proper subfield. Hence there are at most
$\tau(s)$ non-generators. Every proper divisor is at most $\lfloor s/2\rfloor$,
so $\tau(s)\leq\sum_{d=1}^{\lfloor s/2\rfloor}2^d<2^{\lfloor s/2\rfloor+1}$.
The conditional chance that the first coordinate of $v_4$ is not a generator
is at most $2^r\tau(s)/q$. Therefore the clean quadruple identifies the
field with conditional probability at least $1-\beta$, where
\begin{equation}\label{eq:beta}
 \beta=2^R\left(\frac{1+q+q^2}{q^3}+\frac{\tau(s)}q\right).
\end{equation}
This bound holds for every earlier basis triple; it makes no independence
assumption about the field coordinate and $A$.

After computing any ordered binary basis of the correct balance space,
an independent uniform $s$-bit combination is a uniform scalar $t\in K$.
Identify $V^*$ with $K^3$ using the nondegenerate coordinatewise trace
pairing. Under this identification $T^*$ multiplies by $t$. Dependence of
$U_*+tU_*$ yields a nonzero pair $(u_0,u_1)\in U_*^2$ with
$u_0+t u_1=0$. If $u_1=0$ there is no root. Otherwise one nonzero
$K$-coordinate of $u_1$ determines at most one $t$. A union bound over
the $2^{2r}-1$ nonzero pairs thus bounds bad translate position by
$(2^{2r}-1)/q$. Joint generation and good position have probability at least
\begin{equation}\label{eq:gamma}
 \gamma=1-\frac{\tau(s)}q-\frac{2^{2R}-1}q.
\end{equation}

If $s\geq2R+5$, then
\[
 2^R\frac{1+q+q^2}{q^3}<2^{R+1-s}\leq\frac1{16},\qquad
 2^R\frac{\tau(s)}q<2^{R+1-\lceil s/2\rceil}\leq\frac14.
\]
Thus $\beta<1/2$. Also $\tau(s)/q<1/4$ and
$(2^{2R}-1)/q<1/32$, so in particular $\gamma\geq1/4$.
Each trial succeeds with probability at least
$2^{-4R}(1-\beta)\gamma\geq2^{-4R-3}$.
All five draws in each trial are fresh, including the scalar draw in trials
that will fail field identification. With $N=8\cdot2^{4R}h$, independence gives
\[
 \Pr(\text{all trials fail})\leq(1-2^{-4R-3})^N
 \leq e^{-h}\leq2^{-h}.
\]
Earlier acceptance of a different verified representation only helps.

\section{Algorithm and exact certificate}\label{sec:algorithm}
For each of the prescribed $N$ trials perform the following operations.
\begin{enumerate}
\item Draw four $n$-bit vectors and an $s$-bit scalar index. Construct their
contractions directly from the supplied coefficients.
\item Solve the simultaneous balance equations on $n^2$ unknown bits.
Reject unless the solution space has dimension $s$ and contains the identity.
\item Use the scalar index to select $T$ in this space. Form
$I,T,\ldots,T^s$. Require $I,\ldots,T^{s-1}$ independent and in the balance
space, and solve $T^s=\sum_{j<s}p_jT^j$. Test irreducibility of
$p(X)=X^s+\sum_{j<s}p_jX^j$. Reject on failure.
\item In $K=\F_2[X]/(p)$, greedily choose three module generators from the
physical binary basis. For each chosen vector append its first $s$ iterates
under $T$. Require the resulting $n$ columns to be independent and invert
this matrix to obtain $f:V\to K^3$.
\item Compute $U=W(D_TC)$, and reject if its dimension exceeds $R$.
Construct the $s$ trace-determinant coefficient columns in these coordinates,
and the $\binom{\dim U}3$ exterior-cube columns on a basis of $U$.
Solve for the coefficients representing the whole tensor $C$.
\item Reject if there is no solution or the compatible coefficient is zero.
Independently verify the complete serialized certificate, and accept only if
that check passes. After $N$ rejected trials, return failure.
\end{enumerate}

The irreducibility test can compute successive squares of $X$ modulo $p$,
requiring $X^{2^s}=X$ and
$\gcd(p,X^{2^k}-X)=1$ for every $1\leq k\leq\lfloor s/2\rfloor$.
For completeness of this test, a reducible degree-$s$ polynomial has an
irreducible factor of degree at most $s/2$; such a factor divides
$X^{2^k}-X$ for its degree $k$. An irreducible degree-$s$ polynomial satisfies
all the required tests. No factoring oracle or table of field elements is used.

The relation and independence tests make $\F_2[T]$ a field of degree $s$.
Each nonzero vector in the quotient by previously selected $K$-lines generates
another $s$-dimensional binary summand; otherwise a nonzero field scalar would
annihilate its nonzero quotient class. Greedy choice therefore constructs a
$K$-basis in three steps when the field tests pass. Direct verification of
the resulting binary basis is retained.

The certificate contains $p$, the $n\times n$ binary coordinate map $f$,
a nonzero $s$-bit $c$, independent support covectors $u_1,\ldots,u_r$ with
$r\leq R$, and $\binom r3$ binary exterior coefficients. Its size is
$O(n^2+nR+R^3+s)$ bits. The verifier rechecks irreducibility, ranges,
independence, and all $\binom n3$ coefficients of \eqref{eq:model}. Both
summands are alternating trilinear by construction, so checking these
coefficients proves equality on every triple of physical vectors. This is
exact certificate soundness; the verifier never trusts a claimed clean event.

Under a good trial, Lemmas~\ref{lem:four} and~\ref{lem:two} supply the correct
field and the minimal error support. The full coefficient system is therefore
consistent. Moreover its compatible coefficient cannot accidentally be zero:
a nonzero trace-determinant tensor has mode image all of $V^*$, by the trace
nondegeneracy argument, whereas a tensor supported on $U$ has mode image
contained in $U$. Since $\dim U\leq R<n$ under the stated field-size condition,
the compatible tensor space intersects the supported error space only in zero.
Thus any solution to the coefficient system has the same nonzero compatible
component in the recovered coordinates. This completes the correctness and
probability proof of Theorem~\ref{thm:main}.

\section{Complete resource accounting}
All field operations use $s$-bit polynomial representatives. Schoolbook
multiplication and reduction use $O(s^2)$ bit operations; trace evaluation
uses $s$ squarings and $O(s^3)$ bit operations. Binary vectors and matrices
are stored explicitly. The following deliberately loose bounds use elementary
elimination, without assuming unit-cost operations on unbounded integers.

Four contractions require $O(n^3)$ bit operations. The balance solve has
$O(n^2)$ rows and $n^2$ variables, so uses $O(n^6)$ bit operations and
$O(n^4)$ bits of storage. Forming matrix powers, checking their spans, and
constructing and inverting a module basis each fit within $O(n^6)$ bit
operations and $O(n^4)$ storage. The successive-squaring and gcd
irreducibility procedure is polynomial and fits within the same envelope.
Computing the defect slices and their image costs $O(n^4)$ with dense
binary matrix arithmetic and elimination.

The coefficient fit has $d=O(n^3)$ rows and
$m=s+\binom r3=O(n^3)$ columns, even on failed trials with $r\leq R\leq n$.
Constructing the compatible columns costs $O(n^7)$ bit operations using
the trace routine just described. Constructing the exterior columns costs
$O(n^6)$; elimination costs $O(dm\min(d,m))=O(n^9)$ and stores
$O(dm)=O(n^6)$ bits. Certificate rechecking reconstructs these columns
and costs no more. All stages can discard one failed candidate before
starting the next. Polynomial index and interpreter representation overheads
are absorbed conservatively by $\widetilde O(n^{12})$ time and
$\widetilde O(n^8)$ working space per trial. These are envelopes, not optimized
or empirically measured scaling laws.

The trial counter uses $O(R+\log(h+1))$ bits; its arithmetic, input validation,
reading and storage of $C$, and serialization of the certificate are included.
Incremental counter work is at most $O(N(R+\log(h+1)))$. Storage has no
factor $N$. The implementation retains a fixed number of rejection counters,
not all previous candidates. Each trial consumes exactly $4n+s$ bits, even
if an early algebraic test fails. Multiplying the envelope by $N$ gives
\eqref{eq:cost}. At $R=O(\log n)$ and polynomial $h$, this is polynomial
in the dense input size. The exponential dependence on $R$ remains explicit.

\section{Comparison with established methods}\label{sec:prior}
Hora's ring forms~\cite{hora} include composing a determinant over a
characteristic-$p$ ring with a linear functional. His results study radicals,
decomposition and equivalence, including characteristic-two examples. Thus
the uncorrupted form used here is established mathematical structure. The
recovery theorem adds an unknown common support and a bounded randomized
constructor with exact verification; it does not originate ring forms.

Computational tensor invariants already use adjoints, nuclei and centroids
to expose underlying scalar structure. The Magma documentation~\cite{magma}
explicitly includes underlying-field recovery and points to
Brooksbank--Wilson~\cite{bw} for cyclic-algebra machinery. Their module
isomorphism theorem concerns isomorphisms over rings and includes an efficient
cyclic-ring case. Our simultaneous balance solve is elementary linear algebra
within that established approach. The theorem here specifies when four
corrupted contractions happen to identify the field, then certifies the full
input identity. We do not assert a general tensor-isomorphism algorithm.

Fouque--Granboulan--Stern~\cite{fgs} exploit differential kernel dimensions
in multivariate cryptanalysis, including internally perturbed
Matsumoto--Imai systems. This is a direct precedent for finding directions
that cancel a hidden perturbation. Their cryptographic input and inversion
task differ from a scalar-valued alternating cubic tensor and its compatible
representation certificate.

Dubois--Granboulan--Stern~\cite{dgs}, Theorem~1, analyzes a random instance
of internally perturbed HFE and relates derivative-kernel distributions to
random linear maps for a large proportion of instances. Their recovery
uses a statistical distinguisher for the perturbation kernel. In contrast,
our randomness belongs to the algorithm and the corruption tensor is
arbitrary under its support promise. This distinction in quantifiers is
essential; their theorem is not a stated worst-case guarantee for our tensor
class. These comparisons delimit the proposed contribution. They are a
focused primary-source screen, not an exhaustive novelty or priority proof.

\section{Reproducibility, formal status, and limitations}
The companion at
\href{https://github.com/Quantyra/corrupted-field-recovery}{Quantyra/corrupted-field-recovery}
contains \path{src/recovery.py} and \path{tests/check_recovery.py}.
Run the latter with Python 3.13; it writes an environment and source-hash
receipt to \path{artifacts/recovery-checks.json}. The tests distinguish
ordinary seeded tapes from fixture-assisted tapes. They recreate 200 exact
rational probability-bound checks, nine four-contraction field identifications,
and 104 two-translate image checks. These counts are saved executable
evidence, not the unsaved earlier conversation checks.

The suite also exercises coordinate changes, independent pointwise certificate
evaluation, altered certificates and tensors, malformed inputs, a smaller
rank-cap refusal, and all 16 trials of an $R=0,h=2$ zero tape. The $s=11,R=3$
scripted case satisfies the sufficient field-size promise but is not an
ordinary-random success-frequency measurement. The $s=3,R=3$ ordinary seeded
positive has covectors in three independent $K$-directions and is outside
that sufficient condition. A first attempted $s=3,R=3$ fixture instead filled
one entire $K$-coordinate; its seeded run exhausted the bound and is retained
as a development failure, not silently reported as a positive test. The
$s=5$ negative control checks that two three-dimensional subspaces of one
five-dimensional $K$-line cannot be independent, for all 32 scalars.

Lean sources and a theorem-to-obligation inventory accompany this draft.
Compiling individual lemmas is not complete certification of this manuscript.
The four-contraction field theorem, finite-field cardinality and sampling
bridge, complete bounded constructor, verifier equivalence and bit-cost
formalization must all be discharged before a complete certification claim.
The Python receipt alone establishes none of these formal obligations.

The origin is R05-08 at historical revision
\path{d2ae8c6fc275423c2276ba82644b8993660ec684} in
\path{Quantyra/quantyra-quantum-complexity}; exact paths are recorded in
\texttt{PROVENANCE.md}. The present theorem isolates cubic tensors. The
broader phase-recovery application needs additional quadratic and pairing
steps and retains its separate $s\geq7R+5$ sufficient condition. No claim
about complete circuit simulation, quantum advantage, counting backends,
or BPP versus BQP follows by replacing that condition with $2R+5$.

\paragraph{Authorship, assistance, and license.}
The author and affiliation follow established Quantyra project metadata.
Codex assisted with proof reconstruction, implementation, tests, formalization,
literature comparison and manuscript preparation. Only internal review and
the specifically recorded checks have occurred; no independent peer review
or journal acceptance is claimed. Manuscript text is licensed CC BY 4.0;
software and Lean sources are licensed Apache 2.0. No DOI is assigned in this
working draft. Release remains pending full validation and Lean certification.

\begin{thebibliography}{9}
\bibitem{hora} J. Hora, Ring forms, \emph{Italian Journal of Pure and Applied
Mathematics} 41 (2019), 628--636.
\href{https://ijpam.uniud.it/online_issue/201941/55-Jan\%20Hora.pdf}{Primary article}.
\bibitem{magma} Computational Algebra Group, \emph{Magma Handbook:
Linear Invariants of Tensors}, accessed 28 October 2026.
\href{https://docs.magma-maths.org/Modules/MultilinearAlgebra/Linear-invariants-of-tensors.html}{Online documentation}.
\bibitem{bw} P. A. Brooksbank and J. B. Wilson, The module isomorphism
problem reconsidered, \emph{Journal of Algebra} 421 (2015), 541--559.
\href{https://doi.org/10.1016/j.jalgebra.2014.09.004}{doi:10.1016/j.jalgebra.2014.09.004}.
\bibitem{fgs} P.-A. Fouque, L. Granboulan and J. Stern, Differential
cryptanalysis for multivariate schemes, \emph{EUROCRYPT 2005}, LNCS 3494,
341--353. \href{https://www.di.ens.fr/~fouque/pub/eurocrypt05.pdf}{Author manuscript}.
\bibitem{dgs} V. Dubois, L. Granboulan and J. Stern, Cryptanalysis of HFE
with internal perturbation, \emph{PKC 2007}, LNCS 4450, 249--265.
\href{https://www.di.ens.fr/~stern/data/St118.pdf}{Author manuscript}.
\end{thebibliography}
\end{document}
